\section*{6 From Unordered to Ordered Multiselection (author source pp.11:13--11:16)}
In this section we construct a circuit for ordered multiselection over a large alphabet from a circuit for unordered multiselection over a slightly larger alphabet.
\begin{originalresult}{Lemma 22}
Let $q=q(n)\in[n]$ and $s=s(n)=\Omega(\log n)$ be integer functions of $n$, let $\Sigma=\bits^s$, let $\widetilde\Sigma=[n]\times\Sigma$, and suppose there is a (log-space uniform) Boolean circuit for $\uSel_{\widetilde\Sigma}^{n/s\to q}$ with size $S$ and depth $D$.
Then there is a (log-space uniform) Boolean circuit for $\Sel_\Sigma^{n\to q}$ with size $S+O(q\cdot s\cdot\log q)$ and depth $D+O(\log n)$.
\end{originalresult}
Combining Lemma 22 with Proposition 20 gives the following corollary.
\begin{originalresult}{Corollary 23}
Let $s=s(n)$ satisfy $s=\Omega(\log^2n)$, let $\Sigma=\bits^s$, and let $q=q(n)\in[n]$. Then there is a log-space uniform Boolean circuit for $\Sel_\Sigma^{n\to q}$ with size $O(n\cdot s+q\cdot s\cdot\log n)$ and depth $O(\log n)$.
\end{originalresult}
\subsection*{6.1 Boolean Circuits for Inner Joins}
The main tool that we use in the proof of Lemma 22 is Boolean circuitry for computing an inner join of two relations (cf. relational databases [12]).
\begin{originalresult}{Definition 24}
If $R\subseteq W\times X$ and $S\subseteq X\times Y$ are relations, the inner join of $R$ and $S$ (which we denote by $R\bowtie S$) is
\[R\bowtie S\stackrel{\rm def}=\{(w,x,y):(w,x)\in R\land(x,y)\in S\}.\]
\end{originalresult}
To our knowledge, prior work on the computational complexity of computing joins has focused on algorithms in the RAM or PRAM models of computation, as opposed to Boolean circuits.
We focus for simplicity on a special case. First, we require that the relation $S$ is a partial function. That is, for every $x\in X$ there is at most one $y\in Y$ such that $(x,y)\in S$. The partial function requirement prevents $|R\bowtie S|$ from being larger than $|R|$. We also require that for every $(w,x)\in R$ there exists some $(x,y)\in S$.
The following proposition says that inner joins in this special case are computable by (uniform) Boolean circuits of logarithmic depth and nearly linear size.
\begin{originalresult}{Proposition 25}
Let $n,m$ be positive integers, let $W$ and $Y$ be sets whose elements are represented by $s$-bit strings, and let $X$ be a finite set whose elements are represented by $k$-bit strings.
There exists a log-space uniform Boolean circuit of size $O(q\cdot(k+s)\cdot\log q)$ and depth $O(k+\log q)$ that takes as input a relation $R\subseteq W\times X$ and a partial function $f\subseteq X\times Y$ with $|R|\le q$, $|f|\le q$, and $\{x:\exists w,(w,x)\in R\}\subseteq\{x:\exists y,(x,y)\in f\}$, and outputs $R\bowtie f$.
Here sets are represented by an arbitrarily ordered listing of their elements, padded with $\bot$ elements as needed to have length $q$.
\end{originalresult}
\begin{proof}
Denote the elements of $\{x:\exists y,(x,y)\in f\}$ by $x_1,\ldots,x_k$, where $x_1<\cdots<x_k$. For $i\in[k]$, let $W_i$ denote the set $\{w:(w,x_i)\in R\}$, and let $n_i$ denote $|W_i|$. With this notation, our desired output is a list of length $q$ whose non-$\bot$ elements in some order are precisely $(w,x_i,f(x_i))_{i\in[k],w\in W_i}$.
We first construct a list with at most $2q$ elements in $(W\cup\{\star\})\times X\times(Y\cup\{\star\})$, namely $\{(x_i,w,\star)\}_{i\in[k],w\in W_i}$ and $\{(x_i,\star,f(x_i))\}_{i\in[k]}$. Such a list is readily obtained by concatenating the listings of $R$ and $f$ (with the natural embeddings of $W\times X$ and $X\times Y$ in $(W\cup\star)\times X\times(Y\cup\star)$).
We sort this list with respect to the partial ordering that defines $(w,x,y)\prec(w',x',y')$ iff $x<x'$ or $x=x'$ and $w=\star$ and $w'\ne\star$. By Lemma 6 this can be done with circuitry of size $O(q\cdot(k+s)\cdot\log q)$ and depth $O(\log q)$. This yields a list $L$ composed of $k$ blocks, the $i$th of which has length $n_i+1$ and is
\[((\star,x_i,f(x_i)),(w_{i,1},x_i,\star),\ldots,(w_{i,n_i},x_i,\star)),\]
where $\{w_{i1},\ldots,w_{i,n_i}\}=W_i$.
With local processing, we obtain two lists $L_f$ and $L_S$ where the $i$th block of $L_f$ is
\[(f(x_i),\underbrace{\bot,\ldots,\bot}_{n_i\text{ times}})\tag{4}\]
and the $i$th block of $L_S$ is
\[(\bot,(w_{i,1},x_i),\ldots,(w_{i,n_i},x_i)).\tag{5}\]
Using Lemma 26 below, we map $L_f$ to a list whose $i$th block is
\[(\underbrace{f(x_i),\ldots,f(x_i)}_{n_i+1\text{ times}}).\tag{6}\]
We then locally combine (6) with (5) to obtain a list whose $i$th block is
\[(\bot,(w_{i,1},x_i,f(x_i)),\ldots,(w_{i,n_i},x_i,f(x_i))).\tag{7}\]
A final sorting step sends the $\bot$ elements to the back of the list, which allows us to conclude by truncating the list to length $n$. This step can also be done with circuitry of size $O(q\cdot(k+s)\cdot\log q)$ and depth $O(\log q)$ by Lemma 6.
\end{proof}
\begin{originalresult}{Lemma 26}
For $n,s\in\ZZ^+$, there is a (logspace-uniform) constant fan-in Boolean circuit of size $O(n\cdot s)$ and depth $O(\log n)$ that takes as input $x\in(\bits^s\cup\{\bot\})^n$ with $x_1\ne\bot$, and outputs $y\in(\bits^s)^n$ such that
\[y_i=\begin{cases}x_i&\text{if }x_i\ne\bot,\\y_{i-1}&\text{otherwise}.\end{cases}\]
Equivalently, $y_i=x_{j_i}$, where $j_i=\max\{j:1\le j\le i\land x_j\ne\bot\}$ (this maximum is guaranteed to be over a non-empty set because $x_1\ne\bot$).
\end{originalresult}
\begin{proof}
Without loss of generality we can assume that $s=1$ because for $s>1$ we can use $s$ copies of the circuit for the $s=1$ case.
Consider the binary operation $\star:\{0,1,\bot\}\times\{0,1,\bot\}\to\{0,1,\bot\}$ defined by
\[\tau\star\upsilon=\begin{cases}\tau&\text{if }\upsilon=\bot,\\\upsilon&\text{if }\upsilon\ne\bot.\end{cases}\]
It is clear that our desired $y$ has $y_1=x_1$ and $y_i=y_{i-1}\star x_i$ for $i>1$.
We now observe that $\star$ is an associative operation. To see this, consider any $\sigma,\tau,\upsilon\in\{0,1,\bot\}$ and consider separately the cases $\upsilon=\bot$ and $\upsilon\ne\bot$. If $\upsilon=\bot$ then $(\sigma\star\tau)\star\upsilon=\sigma\star\tau=\sigma\star(\tau\star\upsilon)$. If $\upsilon\ne\bot$ then $(\sigma\star\tau)\star\upsilon=\upsilon=\sigma\star\upsilon=\sigma\star(\tau\star\upsilon)$.
Now we use a classic result of Ladner and Fischer [20] that for any associative operation $\star:\Sigma\times\Sigma\to\Sigma$ and $n\in\ZZ^+$, there exists a (log-space uniform) circuit of size $O(n)$ and depth $O(\log n)$ (with only $\star$-gates) that computes all $\star$-prefix products. That is, the circuit takes as input $x\in\Sigma^n$ and outputs $y$ such that $y_i=x_1\star\cdots\star x_i$.
In our case, $\star$ is computable by a Boolean circuit of constant size, so replacing $\star$-gates by such a circuit finishes the proof.
\end{proof}
\subsection*{6.2 A Circuit for Ordered Multiselection}
We are now ready to prove Lemma 22.
\begin{proof}[Proof of Lemma 22]
Let $\widetilde C$ be a circuit for $\uSel_{\widetilde\Sigma}^{n\to q}$.
\paragraph{The Circuit.} Our circuit for $\Sel_\Sigma^{n\to q}$ takes as input $(x,i)\in\Sigma^n\times[n]^q$ and consists of the following stages:
\begin{enumerate}
\item Compute $\widetilde x\in\widetilde\Sigma^n$, defined so that $\widetilde x_i=(i,x_i)$ for $i\in[n]$, and compute a list representation of
\[I=\{(1,i_1),\ldots,(q,i_q)\}.\]
\item Compute $\widetilde y\leftarrow\widetilde C(\widetilde x,i)$ to obtain $\widetilde y\in(\widetilde\Sigma\cup\{\bot\})^q$ such that
\[\{\widetilde y_k:\widetilde y_k\ne\bot\}_{k\in[q]}=\{\widetilde x_{i_k}\}_{k\in[q]}=\{(i_k,x_{i_k})\}_{k\in[q]}.\]
In particular, each $\widetilde y_k$ that is not equal to $\bot$ has the form $(i,x_i)$ for some $i\in[n]$.
\item Deduplicate $\widetilde y$ to obtain a list representation of the partial function
\[Y=\{(i_k,x_{i_k})\}_{k\in[q]}.\]
This involves sorting the non-$\bot$ symbols $(i_k,x_{i_k})$ of $\widetilde y$ in order of increasing $i_k$ (this can be done with circuitry of size $O(n\cdot s\cdot\log n)$ and depth $O(\log n)$ by Lemma 6).
\item Apply the circuit given by Proposition 25 to $I$ and $Y$ to compute a list representation of
\[I\bowtie Y=\{(k,i_k,x_{i_k})\}_{k\in[q]}.\]
\item Sort the list representation of $I\bowtie Y$ in order of increasing $k$, and read off the desired output $(x_{i_1},\ldots,x_{i_k})$.
\end{enumerate}
\paragraph{Size and Depth.} The ``computation'' of $\widetilde x$ from $x$ and of $I$ from $i$ is just adding constant values, so it can certainly be done by a (log-space uniform) Boolean circuit of size $O(n\cdot(s+\log n))$ and depth $O(1)$.
By assumption, $\widetilde C$ is a circuit of size $S$ and depth $D$.
Via sorting (Lemma 6), one can deduplicate $\widetilde y$ with circuitry of size $O(q\cdot(s+\log n)\cdot\log q)=O(q\cdot s\cdot\log q)$ and depth $O(\log n)$.
By Proposition 25, the computation of $I\bowtie Y$ can also be done with circuitry of size $O(q\cdot(s+\log n)\cdot\log q)=O(q\cdot s\cdot\log n)$ and depth $O(\log n)$.
Sorting the elements of $I\bowtie Y$ again takes log-space uniform circuitry of size $O(q\cdot s\cdot\log q)$ and depth $O(\log n)$ by Lemma 6.
In total our constructed circuit has size $S+O(q\cdot s\cdot\log q)$ and depth $D+O(\log n)$.
\end{proof}
